% GENERATED FILE. Edit canonical lessons, then run npm run generate:books.
% canonical-source-hash: fnv1a64:343033b2561616cf
% canonical-lessons: PT-C531-consertar, PT-C531-pendurar, PT-C531-martelar, PT-C531-berbequim, PT-C531-alicate

\chapter{Pequenos arranjos}
\label{ch:pt-pequenos-arranjos}
\hlchaptermodality{531}

\begin{chapteropening}
\textbf{By the end of this chapter:} \emph{I can talk about small repairs at home.}

Complete the last lesson of chapter 531: I can talk about small repairs at home.
\end{chapteropening}

\section[consertar]{consertar — to repair, to fix}

\label{lesson:PT-C531-consertar}



\begin{quote}
Before the new one: say the Portuguese for pantry, then the Portuguese for damp.
\end{quote}

\subsection*{You'll want to know: consertar}
\textbf{consertar} — \textquotedblleft{}to repair, to fix\textquotedblright{}.

\begin{quote}
O meu pai consertou a cadeira partida.
\end{quote}

\begin{quote}
\emph{My father fixed the broken chair.}
\end{quote}

Not \emph{concertar}, which means to arrange or to agree.

\begin{cousinweb}
\emph{consertar} is usually traced to Latin \emph{cōnsertus} \textquotedblleft{}joined together\textquotedblright{}, from \emph{cōnserere} \textquotedblleft{}to join\textquotedblright{} (Lewis \& Short).
\end{cousinweb}

\begin{grammarlens}[title={What you've built}]
One more word for telling, explaining and describing at B1.
\end{grammarlens}

\subsection*{Your turn}
\begin{itemize}
\raggedright
  \item \emph{Say it:} \emph{consertar}
  \item \emph{Say it:} \emph{consertar}, once more
  \item \emph{Say it:} a sentence of your own with \emph{consertar}
  \item \emph{Recall:} say the Portuguese for pantry, then the Portuguese for damp, then say \emph{consertar} again
\end{itemize}

\begin{tcolorbox}[breakable,colback=teal!4,colframe=teal!35!black,title={Before you move on}]
What does consertar mean? (\textquotedblleft{}To repair, to fix\textquotedblright{}.) Say it once more.
\end{tcolorbox}
\section[pendurar]{pendurar — to hang (up)}

\label{lesson:PT-C531-pendurar}



\begin{quote}
Before the new one: say the Portuguese for damp, then the Portuguese for to repair.
\end{quote}

\subsection*{You'll want to know: pendurar}
\textbf{pendurar} — \textquotedblleft{}to hang (up)\textquotedblright{}.

\begin{quote}
Pendurei o quadro na parede da sala.
\end{quote}

\begin{quote}
\emph{I hung the picture on the living-room wall.}
\end{quote}

\emph{Pendurar o casaco} — to hang up your coat.

\begin{cousinweb}
\emph{pendurar} is already found in Old Portuguese (EtymDB) and is usually traced to Latin \emph{pendēre} \textquotedblleft{}to hang\textquotedblright{} (Lewis \& Short).
\end{cousinweb}

\begin{grammarlens}[title={What you've built}]
One more word for telling, explaining and describing at B1.
\end{grammarlens}

\subsection*{Your turn}
\begin{itemize}
\raggedright
  \item \emph{Say it:} \emph{pendurar}
  \item \emph{Say it:} \emph{pendurar}, once more
  \item \emph{Say it:} a sentence of your own with \emph{pendurar}
  \item \emph{Recall:} say the Portuguese for damp, then the Portuguese for to repair, then say \emph{pendurar} again
\end{itemize}

\begin{tcolorbox}[breakable,colback=teal!4,colframe=teal!35!black,title={Before you move on}]
What does pendurar mean? (\textquotedblleft{}To hang (up)\textquotedblright{}.) Say it once more.
\end{tcolorbox}
\section[martelar]{martelar — to hammer}

\label{lesson:PT-C531-martelar}



\begin{quote}
Before the new one: say the Portuguese for to repair, then the Portuguese for to hang.
\end{quote}

\subsection*{You'll want to know: martelar}
\textbf{martelar} — \textquotedblleft{}to hammer\textquotedblright{}.

\begin{quote}
O vizinho martela a parede todas as manhãs.
\end{quote}

\begin{quote}
\emph{The neighbour hammers at the wall every morning.}
\end{quote}

Built on \emph{o martelo}.

\begin{cousinweb}
\emph{martelar} is built on \emph{martelo} \textquotedblleft{}hammer\textquotedblright{}, which goes back to Late Latin \emph{martellus} (EtymDB).
\end{cousinweb}

\begin{grammarlens}[title={What you've built}]
One more word for telling, explaining and describing at B1.
\end{grammarlens}

\subsection*{Your turn}
\begin{itemize}
\raggedright
  \item \emph{Say it:} \emph{martelar}
  \item \emph{Say it:} \emph{martelar}, once more
  \item \emph{Say it:} a sentence of your own with \emph{martelar}
  \item \emph{Recall:} say the Portuguese for to repair, then the Portuguese for to hang, then say \emph{martelar} again
\end{itemize}

\begin{tcolorbox}[breakable,colback=teal!4,colframe=teal!35!black,title={Before you move on}]
What does martelar mean? (\textquotedblleft{}To hammer\textquotedblright{}.) Say it once more.
\end{tcolorbox}
\section[o berbequim]{o berbequim — (electric) drill}

\label{lesson:PT-C531-berbequim}



\begin{quote}
Before the new one: say the Portuguese for to hang, then the Portuguese for to hammer.
\end{quote}

\subsection*{You'll want to know: o berbequim}
\textbf{o berbequim} — \textquotedblleft{}(electric) drill\textquotedblright{}.

\begin{quote}
Fez um furo na parede com o berbequim.
\end{quote}

\begin{quote}
\emph{He drilled a hole in the wall with the drill.}
\end{quote}

In Brazil: \emph{a furadeira}.

\begin{cousinweb}
\emph{berbequim} is usually traced to French \emph{vilebrequin}, a hand drill; the origin of the French word is not settled.
\end{cousinweb}

\begin{grammarlens}[title={What you've built}]
One more word for telling, explaining and describing at B1.
\end{grammarlens}

\subsection*{Your turn}
\begin{itemize}
\raggedright
  \item \emph{Say it:} \emph{o berbequim}
  \item \emph{Say it:} \emph{o berbequim}, once more
  \item \emph{Say it:} a sentence of your own with \emph{o berbequim}
  \item \emph{Recall:} say the Portuguese for to hang, then the Portuguese for to hammer, then say \emph{o berbequim} again
\end{itemize}

\begin{tcolorbox}[breakable,colback=teal!4,colframe=teal!35!black,title={Before you move on}]
What does o berbequim mean? (\textquotedblleft{}(electric) drill\textquotedblright{}.) Say it once more.
\end{tcolorbox}
\section[o alicate]{o alicate — pliers}

\label{lesson:PT-C531-alicate}



\begin{quote}
Before the new one: say the Portuguese for to hammer, then the Portuguese for (electric) drill.
\end{quote}

\subsection*{You'll want to know: o alicate}
\textbf{o alicate} — \textquotedblleft{}pliers\textquotedblright{}.

\begin{quote}
Usou o alicate para tirar o prego.
\end{quote}

\begin{quote}
\emph{He used the pliers to pull out the nail.}
\end{quote}

One tool in Portuguese: \emph{um alicate}.

\begin{cousinweb}
\emph{alicate} comes from Arabic \emph{laqqāṭ} \textquotedblleft{}pincers\textquotedblright{} (EtymDB), with the Arabic article \emph{al-}.
\end{cousinweb}

\begin{grammarlens}[title={What you've built}]
One more word for telling, explaining and describing at B1.
\end{grammarlens}

\subsection*{Your turn}
\begin{itemize}
\raggedright
  \item \emph{Say it:} \emph{o alicate}
  \item \emph{Say it:} \emph{o alicate}, once more
  \item \emph{Say it:} a sentence of your own with \emph{o alicate}
  \item \emph{Recall:} say the Portuguese for to hammer, then the Portuguese for (electric) drill, then say \emph{o alicate} again
\end{itemize}

\begin{tcolorbox}[breakable,colback=teal!4,colframe=teal!35!black,title={Before you move on}]
What does o alicate mean? (\textquotedblleft{}Pliers\textquotedblright{}.) Say it once more.
\end{tcolorbox}
